\section{Biotecnologie}

% ---------------------------------------------------------------------------
\subsection{Microscopia e preparazione dei campioni}

Strumenti di indagine: microscopio ottico e microscopio elettronico (a trasmissione, a scansione).

\rubrica{Preparazione dei campioni per il microscopio elettronico}
$\text{fissazione} \rightarrow \text{disidratazione (etanolo 70--100\%)} \rightarrow \text{inclusione in resina} \rightarrow \text{sezionamento} \rightarrow \text{colorazione}$.
\begin{itemize}
  \item fissazione (es.\ glutaraldeide, $OsO_4$) e disidratazione in etanolo a concentrazione crescente;
  \item infiltrazione e inclusione in resina, polimerizzata in un blocco solido;
  \item sezionamento con \textbf{ultramicrotomo} in sezioni di $\sim$100 nm;
  \item colorazione delle sezioni con metalli pesanti.
\end{itemize}

% ---------------------------------------------------------------------------
\subsection{Centrifugazione in gradiente di densità}

Purificazione di frazioni subcellulari: in un gradiente di densità (saccarosio) i vari organelli sedimentano finché raggiungono il livello a densità uguale alla loro, dove formano bande.
\begin{itemize}
  \item lisosomi: 1,12 g/ml;
  \item mitocondri: 1,18 g/ml;
  \item perossisomi: 1,23 g/ml.
\end{itemize}

\rubrica{Sedimentazione degli acidi nucleici}
\begin{itemize}
  \item separazione del DNA in base alle \textbf{dimensioni} (velocità di sedimentazione);
  \item separazione in base alla \textbf{composizione in basi} (centrifugazione all'equilibrio).
\end{itemize}

\rubrica{Gradiente di cloruro di cesio (CsCl)}
Centrifugazione all'equilibrio: separa il DNA in base alla densità.
\begin{itemize}
  \item DNA ricco in AT: $\sim$1,65 g/ml; DNA ricco in GC: $\sim$1,75 g/ml;
  \item separa il DNA plasmidico dal DNA cromosomico batterico (due bande);
  \item il DNA è reso visibile con \textbf{bromuro di etidio}, fluorescente ai raggi ultravioletti.
\end{itemize}

% ---------------------------------------------------------------------------
\subsection{Enzimi di restrizione}

\textbf{Endonucleasi di restrizione}: tagliano il DNA a livello di sequenze palindromiche di riconoscimento, producendo estremità complementari dette \textit{estremità coesive} (``sticky ends'').
\begin{itemize}
  \item esempi: EcoRI (da \textit{E.\ coli}), HindII e HindIII (da \textit{Haemophilus influenzae}), HpaII, AluI;
  \item la \textbf{metilazione} della sequenza di riconoscimento protegge il DNA dal taglio.
\end{itemize}

% ---------------------------------------------------------------------------
\subsection{Clonaggio e DNA ricombinante}

\rubrica{Vettori plasmidici}
Plasmidi (es.\ pUC118, pUC119) che contengono:
\begin{itemize}
  \item una \textbf{regione di policlonaggio} con i siti di taglio per gli enzimi di restrizione;
  \item il gene \textit{lacZ};
  \item un'origine di replicazione;
  \item un gene di resistenza all'ampicillina.
\end{itemize}

\rubrica{Formazione di DNA ricombinante}
Lo stesso enzima di restrizione taglia sia il plasmide sia il DNA da inserire $\rightarrow$ i frammenti hanno estremità adesive complementari $\rightarrow$ il frammento estraneo si inserisce nel plasmide e le interruzioni vengono saldate dalla \textbf{DNA ligasi}.

% ---------------------------------------------------------------------------
\subsection{\texorpdfstring{PCR (\textit{Polymerase Chain Reaction})}{PCR (Polymerase Chain Reaction)}}

\textbf{Scopo}: amplificare in provetta una sequenza bersaglio, cioè produrne un grande numero di copie, senza clonaggio.

\rubrica{Componenti della miscela di reazione}
\begin{enumerate}
  \item il DNA contenente la sequenza bersaglio da amplificare;
  \item una coppia di primer, ciascuno complementare a una delle estremità della sequenza bersaglio;
  \item i quattro nucleotidi trifosfato (dATP, dCTP, dGTP, dTTP);
  \item DNA polimerasi \textbf{termostabile} (isolata da un microrganismo termofilo, perché la PCR usa alte temperature in alcuni passaggi).
\end{enumerate}
La reazione è ciclica: a ogni ciclo il numero di copie raddoppia ($2 \rightarrow 4 \rightarrow 8 \rightarrow \dots$).

% ---------------------------------------------------------------------------
\subsection{Elettroforesi su gel di agarosio}

Separa molecole di DNA, RNA o proteine in base alla dimensione, alla carica o ad altre proprietà, applicando un campo elettrico a un gel.
\begin{enumerate}
  \item si prepara un gel di agarosio con pozzetti e lo si immerge in un tampone tra due elettrodi;
  \item si caricano nei pozzetti i campioni (es.\ prodotti di PCR) e un \textit{marker} (frammenti di dimensione nota);
  \item applicando la corrente, i frammenti (carichi negativamente) migrano verso il polo positivo; i più corti migrano più velocemente;
  \item i frammenti della stessa lunghezza formano bande, visualizzate con un colorante che lega il DNA e diventa fluorescente ai raggi UV.
\end{enumerate}

% ---------------------------------------------------------------------------
\subsection{Sequenziamento del DNA}

Metodo con terminatori di catena: si preparano quattro reazioni di sintesi, ciascuna contenente DNA polimerasi, primer, i quattro dNTP e uno dei quattro \textbf{didesossiribonucleosidi trifosfato} (ddNTP), marcato.
\begin{itemize}
  \item l'incorporazione di un ddNTP interrompe la sintesi $\rightarrow$ si generano frammenti di lunghezza diversa;
  \item i prodotti sono separati per elettroforesi (gel di poliacrilammide o elettroforesi capillare) in base alla grandezza;
  \item la lettura fornisce la sequenza; nei sistemi automatici i ddNTP sono marcati con coloranti fluorescenti letti da un laser e un rilevatore collegato a un computer.
\end{itemize}

% ---------------------------------------------------------------------------
\subsection{Progetto Genoma Umano}

\begin{itemize}
  \item tempo previsto: 1990--2005; tempo effettivamente impiegato: 1990--2002;
  \item sequenziato tutto il genoma umano ($3 \cdot 10^9$ bp);
  \item solo il 3\% è costituito da DNA codificante;
  \item il restante 97\% è detto intergenico o non codificante.
\end{itemize}

% ---------------------------------------------------------------------------
\subsection{Organismi transgenici e clonazione}

\rubrica{Topi transgenici}
Scopo: produrre un animale transgenico che possa trasmettere il transgene alla progenie.
\begin{enumerate}
  \item si introduce il gene desiderato nelle cellule della linea germinale (prelevate da un embrione) tramite iniezione o elettroporazione;
  \item si clona la cellula che ha incorporato il transgene, ottenendo una coltura pura di cellule transgeniche;
  \item si iniettano le cellule transgeniche in embrioni precoci (blastocisti);
  \item si impiantano gli embrioni in madri adottive e si incrocia tra loro la progenie $\rightarrow$ animale ingegnerizzato geneticamente.
\end{enumerate}
Esempio: topo gigante ottenuto introducendo il gene di ratto per l'ormone della crescita.

\rubrica{Clonazione}
Febbraio 1997: clonazione della pecora \textbf{Dolly}, primo mammifero clonato a partire da una cellula di un adulto.

% ---------------------------------------------------------------------------
\subsection{Cellule staminali}

\begin{itemize}
  \item \textbf{cellule totipotenti} --- cellule staminali embrionali (CSE): in grado di produrre qualsiasi cellula dell'organismo;
  \item \textbf{cellule pluripotenti o multipotenti} --- cellule staminali dell'adulto (CSA): possono dare origine a un certo numero di tessuti (es.\ le staminali ematopoietiche del midollo osseo generano le cellule del sangue).
\end{itemize}
